% fm-03-prior-art

\section{Prior Art: Golden-Ratio Tradition in Fine-Structure Constant Research}

This chapter situates the Trinity--Pellis Bridge inside the broader $\varphi$-structured fine-structure-constant (FSC) tradition. Three independent lines of prior art are essential context for any reader evaluating Strand III (Pellis Hierarchical Expansion): Heyrovska's golden-angle interpretation, Sherbon's Golden Ratio Geometry framework, and Olsen's Tier-D $\varphi$-cosmology cross-reference. None of these prior workers are co-authors of this compendium; their work is cited as independent prior art that the Pellis Hierarchical Expansion methodology builds upon and that the present bridge does not duplicate.

\subsection{1. Heyrovska --- Golden-Angle Interpretation of $\alpha^{-1}$}

Raji Heyrovska (Czech Academy of Sciences, Heyrovský Institute of Physical Chemistry) was among the first to interpret the fine-structure constant in terms of the golden ratio $\varphi$ = (1+$\sqrt{5}$)/2. Her work proposes that the inverse fine-structure constant $\alpha^{-1}$ is related to the golden angle 36$0^{\circ}$/$\varphi^{2}$ ($\approx$ 137.507$7^{\circ}$) and to the Bohr-radius-based ground-state geometry of hydrogen. The relevant identities are:

- 36$0^{\circ}$/$\varphi^{2}$ $\approx$ 137.5077$6^{\circ}$
- $\alpha^{-1}$ (CODATA 2018) = 137.035999084(21)
- $\Delta$ $\approx$ 0.47$2^{\circ}$, which Heyrovska's framework attributes to a small relativistic/QED correction term in a $\varphi$-additive expansion.

Heyrovska's interpretation directly anticipates the Pellis additive expansion \_\_C0\_\_ (Pellis viXra 2110.0117). In this compendium, Heyrovska's identity is treated as \textbf{foundational prior art for Strand III} and is the historical entry point of golden-ratio FSC research. Source: Heyrovska, R., "The Golden Ratio in the Creations of Nature Arises in the Architecture of Atoms and Ions", various publications collected on her institutional page.

\subsection{2. Sherbon --- Golden Ratio Geometry and the FSC}

Michael A. Sherbon's "Wolfgang Pauli and the Fine-Structure Constant" and "Fundamental Nature of the Fine-Structure Constant" (HAL hal-02341850, https://hal.science/hal-02341850/document) develop a Golden Ratio Geometry framework in which $\alpha$ emerges from $\varphi$-based regular geometric constructions (icosahedral and dodecahedral symmetries, golden-ratio rectangles inscribed in spherical caps). Sherbon explicitly connects:

- The Pauli "magic number" 137 to $\varphi$-based polyhedral angle defects.
- The dodecahedral / icosahedral symmetry group $A_{5}$ (order 60) to the appearance of 36$0^{\circ}$ in golden-ratio FSC formulas.
- Pythagorean / Platonic harmonic ratios to the additive $\varphi^{-k}$ structure later formalised by Pellis.

Sherbon's framework is the second pillar of Strand-III prior art. The Pellis Hierarchical Expansion (this compendium, Strand III) can be read as a fully numerical realisation of Sherbon's geometric intuition.

\subsection{3. Olsen --- Tier-D $\varphi$-Cosmology Cross-Reference}

Scott Olsen ("The Golden Section: Nature's Greatest Secret", Wooden Books) provides the third pillar: a Tier-D cosmological cross-reference in which $\varphi$ appears as an organising ratio across scales, from atomic to galactic. Olsen's work is cited in this compendium \textbf{only as cross-reference} --- it does not enter the falsification ledger directly.

\subsection{4. Coldea et al. --- Empirical E8 Anchor}

The empirical anchor of Strand III is Coldea et al., "Quantum Criticality in an Ising Chain: Experimental Evidence for Emergent E8 Symmetry", Science 327 (2010) 177 [DOI: 10.1126/science.1180085], with the 2023 update Wu et al., "E8 Symmetry in the Quantum Critical Ising Chain CoN$b_{2}$$O_{6}$ Confirmed at Higher Resolution", Phys. Rev. B 108, L060403 (2023) [DOI: 10.1103/PhysRevB.108.L060403]. The Coldea experiment measured the lowest two excitation masses in CoN$b_{2}$$O_{6}$ and found their ratio $m_{2}$/$m_{1}$ = 2 cos($\pi$/30) $\approx$ $\varphi$, providing the first experimental observation of an E8 mass ratio in a condensed-matter system. The 2023 update confirms the original ratio at higher resolution and extends the measurement to three masses, with all ratios within E8 predictions.

The theoretical backbone connecting Coldea's experiment to the $\varphi$-additive structure is Fring \& Korff, "Affine Toda field theories related to Coxeter groups of non-crystallographic type", arXiv:hep-th/0506226 (Nucl. Phys. B 729 (2005) 361) [DOI: 10.1016/j.nuclphysb.2005.08.044]. Fring-Korff show that $\varphi$ appears as the smallest mass ratio in non-crystallographic affine Toda field theories ($H_{2}$, $H_{3}$, $H_{4}$), which is the rigorous mathematical home of the Coldea ratio.

\subsection{5. How the Bridge Differs}

Trinity-Pellis Bridge differs from the prior art in three concrete ways:

1. \textbf{Two-tier filtration.} Pellis publishes additive $\varphi^{-k}$ formulas. The Bridge formalises the filtration \_\_C1\_\_ and proves that Pellis expansions live in F\_{$-$5} (golden-angle) or F\_{$-$2} (Archimedes), giving an MDL-compatible cost.
2. \textbf{Silicon anchor.} Heyrovska, Sherbon, and Pellis present FSC formulas; none anchor them in physical hardware. Strand II (TTSKY26b 0x47C0) attaches a verifiable, reset-time silicon witness to the Trinity identity $\varphi^{2}$ + $\varphi^{-2}$ = 3.
3. \textbf{Falsification ledger.} Coldea, Heyrovska, and Sherbon do not provide a pre-registered falsification protocol. The Bridge supplies one (Strand I + Catalog42 freeze protocol).

\subsection{6. Reading order for Strand-III evaluators}

A reader entering Strand III should consult, in order: Heyrovska (golden-angle origin) $\rightarrow$ Sherbon (geometric framework) $\rightarrow$ Pellis viXra 2110.0117 (numerical expansion) $\rightarrow$ Coldea Science 327 (empirical E8 anchor) $\rightarrow$ Wu et al. PhysRevB.108 (2023 update) $\rightarrow$ Fring-Korff (theoretical backbone) $\rightarrow$ this compendium (Bridge construction).

\subsection{Citations}

- Heyrovska, R., golden-angle FSC interpretation --- institutional publications, Heyrovský Institute.
- Sherbon, M. A., "Wolfgang Pauli and the Fine-Structure Constant" / "Fundamental Nature of the Fine-Structure Constant", HAL hal-02341850, https://hal.science/hal-02341850/document.
- Olsen, S., "The Golden Section: Nature's Greatest Secret", Wooden Books, 2006.
- Coldea, R. et al., Science 327 (2010) 177, DOI: 10.1126/science.1180085.
- Wu, J. et al., Phys. Rev. B 108, L060403 (2023), DOI: 10.1103/PhysRevB.108.L060403.
- Fring, A.; Korff, C., Nucl. Phys. B 729 (2005) 361, arXiv:hep-th/0506226, DOI: 10.1016/j.nuclphysb.2005.08.044.
- Pellis, S., "The Fine-Structure Constant and the Golden Ratio", viXra 2110.0117, https://vixra.org/pdf/2110.0117v4.pdf.
